% Manual de usuario de LabeSAT.
% Fuente canónica en LaTeX (ADR-0006). La compila CI: trabajo «documentacion».
% Local:  latexmk -pdf -cd docs/manual/manual-usuario.tex
\documentclass[11pt,a4paper]{article}

\usepackage[utf8]{inputenc}
\usepackage[T1]{fontenc}
\usepackage[spanish,es-noquoting]{babel}
\usepackage{lmodern}
\usepackage{microtype}
\usepackage[margin=2.5cm]{geometry}
\usepackage{booktabs}
\usepackage{tabularx}
\usepackage{xcolor}
\usepackage{listings}
\usepackage[hidelinks]{hyperref}

% Macros compartidas por todos los documentos LaTeX del proyecto (ADR-0006).
\newcommand{\labesat}{\textsc{LabeSAT}}
\newcommand{\kissat}{\textsc{Kissat}}
\newcommand{\satsuma}{\textsc{satsuma}}
% Versión del documento: la fija CI con -usepretex, o se queda en «borrador».
\providecommand{\versiondoc}{borrador}


\lstset{
  basicstyle=\ttfamily\small,
  columns=fullflexible,
  keepspaces=true,
  frame=single,
  framerule=0.3pt,
  rulecolor=\color{black!30},
  backgroundcolor=\color{black!3},
  showstringspaces=false,
  breaklines=true,
  % listings no entiende UTF-8 multibyte: las tildes de los comentarios se
  % mapean a mano (si no, pdflatex aborta dentro de los bloques de código).
  literate={á}{{\'a}}1 {é}{{\'e}}1 {í}{{\'i}}1 {ó}{{\'o}}1 {ú}{{\'u}}1
           {Á}{{\'A}}1 {É}{{\'E}}1 {Í}{{\'I}}1 {Ó}{{\'O}}1 {Ú}{{\'U}}1
           {ñ}{{\~n}}1 {Ñ}{{\~N}}1 {ü}{{\"u}}1 {→}{{$\rightarrow$}}1,
}

\title{\labesat{} --- Manual de usuario}
\author{Abel Ponce}
\date{Versión del documento: \versiondoc}

\begin{document}
\maketitle

\begin{abstract}
\noindent
\labesat{} es un solver SAT para fórmulas en forma normal conjuntiva (formato
DIMACS). Se basa en Kissat~4.0.4 y añade un preprocesado de \emph{ruptura de
simetrías} (satsuma) cuyo resultado se puede verificar. Este manual explica
cómo instalarlo, cómo usarlo, cómo verificar sus respuestas y cómo reproducir
los experimentos del proyecto.
\end{abstract}

\tableofcontents

%======================================================================
\section{Qué es \labesat{}}

\labesat{} decide si una fórmula proposicional en CNF es satisfacible:
\begin{itemize}
  \item si lo es, imprime un \emph{modelo}, es decir, una asignación que la
        satisface;
  \item si no lo es, puede escribir una \emph{prueba} de insatisfacibilidad que
        un verificador independiente comprueba.
\end{itemize}

Consta de dos programas que se ejecutan en tubería:
\begin{enumerate}
  \item \textbf{satsuma} (M.~Anders, licencia MIT) busca simetrías de la
        fórmula y añade restricciones que las rompen. Escribe la fórmula
        simplificada y el comienzo de la prueba en formato SR
        (\emph{substitution redundancy}).
  \item \textbf{kissat} (A.~Biere, licencia MIT, con modificaciones de
        \labesat{}) resuelve la fórmula simplificada y \emph{continúa} la misma
        prueba.
\end{enumerate}

El guion \texttt{solver/labesat} coordina los dos. Si satsuma falla o tarda
demasiado, resuelve la fórmula original solo con kissat.

%======================================================================
\section{Instalación}

\subsection{Requisitos}
\begin{itemize}
  \item Linux o macOS, con compilador de C (gcc o clang) y \texttt{make}.
  \item Para la ruptura de simetrías: compilador de C++20, CMake~$\geq$~3.5 y
        \texttt{git}.
  \item Python~3.10 o superior para los scripts de experimentos
        (\texttt{pip install -r requirements.txt}).
  \item \texttt{xz}, \texttt{gzip} o \texttt{bzip2} si las instancias vienen
        comprimidas.
\end{itemize}

\subsection{Compilación}
\begin{lstlisting}[language=bash]
git clone https://github.com/abelponce03/SAT-SOLVER.git
cd SAT-SOLVER
./scripts/build.sh          # compila solver/kissat/build/kissat
./scripts/get_tools.sh      # satsuma (MIT), dejavu, dsr-trim y drat-trim en tools/
./scripts/smoke_test.sh     # comprobación rápida de que todo funciona
\end{lstlisting}

\texttt{get\_tools.sh} descarga cada herramienta fijada a un commit concreto,
para que cualquier resultado sea reproducible. satsuma se compila \emph{sin}
la biblioteca cliquer (GPL), de modo que todo lo que se distribuye es MIT.

\subsection{Configuración de competición}
\texttt{./scripts/build.sh --competition} compila la configuración que se
entrega a la SAT Competition. No admite opciones en tiempo de ejecución, salvo
las de la aplicación (\texttt{-n}, \texttt{--time}, \texttt{--append-proof},
\ldots).

%======================================================================
\section{Uso}

\subsection{Línea de órdenes}
\begin{lstlisting}
solver/labesat [--no-symmetry] [opciones de kissat] <cnf> [<prueba>]
\end{lstlisting}

\begin{description}
  \item[\texttt{<cnf>}] Fórmula en formato DIMACS, sin comprimir o comprimida
        (\texttt{.xz}, \texttt{.gz}, \texttt{.bz2}, \texttt{.lzma}).
  \item[\texttt{<prueba>}] Fichero donde se escribe la prueba si la fórmula es
        insatisfacible. Opcional.
  \item[\texttt{--no-symmetry}] Desactiva satsuma y resuelve solo con kissat.
  \item[\texttt{--time=N}] Límite de \textbf{tiempo total} en segundos. El
        tiempo que consume satsuma se descuenta del de kissat.
  \item[otras opciones] Cualquier opción que empiece por \texttt{-} se pasa a
        kissat; por ejemplo, \texttt{-n} (no imprimir el modelo),
        \texttt{-q} o \texttt{--seed=S}. Lista completa: \texttt{kissat -h}.
\end{description}

\subsection{Ejemplos}
\begin{lstlisting}[language=bash]
# Resolver una fórmula y ver el modelo, si lo hay
solver/labesat formula.cnf

# Resolver con límite de 5000 s y escribir la prueba
solver/labesat --time=5000 formula.cnf.xz prueba.out

# Solo kissat, sin ruptura de simetrías
solver/labesat --no-symmetry formula.cnf
\end{lstlisting}

\subsection{Salida y códigos de salida}
La salida sigue el formato de la SAT Competition:
\begin{itemize}
  \item las líneas que empiezan por \texttt{c} son comentarios;
  \item exactamente una línea \texttt{s} con la respuesta;
  \item si la fórmula es satisfacible, líneas \texttt{v} con el modelo,
        terminadas en \texttt{0}.
\end{itemize}
Las líneas \texttt{c [labesat]} informan de qué hizo la tubería.

\begin{center}
\begin{tabular}{@{}lll@{}}
\toprule
Código & Línea de solución & Significado \\
\midrule
10 & \texttt{s SATISFIABLE}   & Hay modelo (líneas \texttt{v}) \\
20 & \texttt{s UNSATISFIABLE} & No hay modelo; prueba en \texttt{<prueba>} \\
0  & \texttt{s UNKNOWN}       & Se agotó el tiempo o se interrumpió \\
\bottomrule
\end{tabular}
\end{center}

\subsection{Variables de entorno}
\begin{center}
\begin{tabularx}{\linewidth}{@{}lXl@{}}
\toprule
Variable & Efecto & Por defecto \\
\midrule
\texttt{LABESAT\_KISSAT} & Binario de kissat & \texttt{solver/kissat/build/kissat} \\
\texttt{LABESAT\_SATSUMA} & Binario de satsuma & \texttt{tools/satsuma} \\
\texttt{LABESAT\_SYMM\_TIMEOUT} & Segundos máximos para satsuma & 60 \\
\texttt{LABESAT\_SYMM\_MAXBYTES} & Tamaño máximo de la CNF para aplicar satsuma & 512\,MiB \\
\texttt{LABESAT\_KISSAT\_ARGS} & Argumentos extra para kissat & (vacío) \\
\bottomrule
\end{tabularx}
\end{center}
Los valores por defecto de los topes son provisionales y se fijarán con el
experimento EXP-007.

%======================================================================
\section{Verificar las respuestas}

Una respuesta solo vale si se puede comprobar de forma independiente.

\paragraph{Satisfacible.} El modelo tiene que satisfacer la fórmula
\emph{original}, no la simplificada:
\begin{lstlisting}[language=bash]
solver/labesat formula.cnf > salida.txt
python3 scripts/verify_model.py --model salida.txt formula.cnf
\end{lstlisting}

\paragraph{Insatisfacible.} La prueba combina el prefijo SR de satsuma y los
pasos DRAT de kissat, así que se verifica con \texttt{dsr-trim}:
\begin{lstlisting}[language=bash]
solver/labesat formula.cnf prueba.out
tools/dsr-trim formula.cnf prueba.out           # debe imprimir: s VERIFIED UNSAT
tools/dsr-trim-sc2026 formula.cnf prueba.out    # versión exacta de la SAT Competition 2026
\end{lstlisting}
Si se usó \texttt{--no-symmetry}, o si satsuma no llegó a actuar, la prueba es
DRAT pura y también la acepta \texttt{tools/drat-trim}.

%======================================================================
\section{Reproducir los experimentos}

Cada experimento está descrito en \texttt{docs/experiments/EXP-NNN-*.md}, que
recoge hipótesis, diseño, criterio de decisión e instrucciones exactas para
reproducirlo. Por ejemplo, el A/B de la ruptura de simetrías (EXP-007):
\begin{lstlisting}[language=bash]
python3 scripts/select_symm_bench.py --download
python3 scripts/run_ab_interleaved.py --solver solver/labesat \
    --guard solver/kissat/build/kissat --guard tools/satsuma \
    --bench bench/symm2026 --out-a A.csv --out-b B.csv \
    --opts-a="--no-symmetry" --timeout 180 --seeds 1
python3 scripts/par2.py A.csv B.csv
\end{lstlisting}

%======================================================================
\section{Solución de problemas}

\begin{description}
  \item[\texttt{labesat: no encuentro kissat}] Compila con
        \texttt{./scripts/build.sh} o indica el binario con
        \texttt{LABESAT\_KISSAT}.
  \item[\texttt{c [labesat] sin simetrías: satsuma no disponible}] Ejecuta
        \texttt{./scripts/get\_tools.sh}. \labesat{} funciona igualmente, pero
        sin ruptura de simetrías.
  \item[\texttt{failed to open and append proof}] \texttt{--append-proof} no
        admite ficheros de prueba comprimidos. Usa un nombre sin
        \texttt{.xz}/\texttt{.gz}.
\end{description}

%======================================================================
\section{Licencias y créditos}

\labesat{} se distribuye bajo licencia MIT. Incorpora Kissat (MIT; A.~Biere,
M.~Fleury, F.~Pollitt) y usa satsuma y dejavu (MIT; M.~Anders). Para verificar
usa dsr-trim (Apache~2.0; C.~R.~Codel) y drat-trim (MIT; M.~Heule,
N.~Wetzler). El detalle está en \texttt{THIRD\_PARTY\_NOTICES.md}.

\end{document}
